\frfilename{mt26.tex}
\versiondate{5.9.03/24.7.04}
\copyrightdate{2004}

\def\diam{\mathop{\text{diam}}}
\def\dist{\mathop{\text{dist}}}

\def\chaptername{Perubahan Variabel dalam Integral}
\newchapter{26}
\def\chaptername{Perubahan variabel dalam integral}

Saya kira kebanyakan mata kuliah kalkulus dasar masih menyediakan waktu
yang cukup banyak untuk berlatih teknik mengintegralkan fungsi-fungsi
standar. Tentulah teknik tunggal yang paling ampuh ialah substitusi:
mengganti $\int g(y)dy$ dengan $\int g(\phi(x))\phi'(x)dx$ untuk suatu
fungsi $\phi$ yang sesuai. Pada tingkat ini, perhatian biasanya dipusatkan
pada keterampilan menebak $\phi$ yang sesuai dan menuliskan rumus-rumus
dengan tepat. Saya tidak akan membahas persoalan semacam itu di sini,
kecuali dalam beberapa kasus khusus yang jarang; dalam buku ini perhatian
saya justru tertuju pada pembenaran proses tersebut. Untuk fungsi satu
variabel, proses ini biasanya dapat dibenarkan dengan merujuk pada Teorema
Dasar Kalkulus, dan untuk setiap kasus tertentu saya biasanya akan terlebih
dahulu menuju \S225 dengan harapan hasil-hasil di sana mencakupnya. Namun,
untuk fungsi dua variabel atau lebih diperlukan beberapa gagasan yang jauh
lebih mendalam.

Saya telah membahas persoalan umum integrasi dengan substitusi dalam ruang
ukur abstrak di \S235. Di sana saya menguraikan syarat-syarat yang menjamin
$\int g(y)dy=\int g(\phi(x))J(x)dx$ untuk suatu fungsi $J$ yang sesuai.
Konteks di sana hampir tidak memberi petunjuk tentang cara menghitung $J$;
paling banter, fungsi tersebut dapat disajikan sebagai turunan
Radon--Nikod\'ym (235M).
Dalam bab ini saya memberikan suatu bentuk teorema dasar untuk kasus ukuran
Lebesgue, dengan $\phi$ suatu fungsi yang sedikit banyak terdiferensialkan
antara ruang-ruang Euklides, dan $J$ suatu `Jacobian', yaitu modulus
determinan turunan $\phi$ (263D). Hal ini dengan sendirinya bergantung pada
penyelidikan yang serius tentang hubungan antara ukuran Lebesgue dan
geometri. Langkah pertama ialah menetapkan suatu bentuk teorema Vitali untuk
ruang berdimensi $r$, beserta teorema-teorema kerapatan berdimensi $r$;
saya mengerjakannya dalam \S261 dengan mengikuti secara dekat rancangan
\S\S221 dan 223 di atas. Kita perlu mengetahui cukup banyak tentang fungsi
terdiferensialkan antara ruang-ruang Euklides, dan ternyata teorinya
berjalin dengan teori fungsi `Lipschitz'; keduanya saya bahas dalam \S262.

Dalam dua bagian berikutnya dari bab ini, saya beralih ke suatu persoalan
terpisah yang ternyata cocok ditangani dengan beberapa teknik yang sama:
penguraian ukuran permukaan pada permukaan (mulus) dalam ruang Euklides,
seperti permukaan kerucut atau bola. Saya kira tidak sulit membentuk intuisi
yang kukuh tentang apa yang dimaksud dengan `luas' permukaan demikian dan
daerah-daerah yang cukup sederhana di dalamnya, dan terdapat dugaan yang
sangat kuat bahwa intuisi ini seharusnya dapat diungkapkan dalam istilah
teori ukuran sebagaimana disajikan dalam buku ini; tetapi menurut saya
rinciannya tidak sederhana. Hal pertama yang perlu diperhatikan ialah bahwa
untuk setiap perhitungan luas suatu daerah $G$ pada permukaan $S$, kita akan
selalu langsung beralih ke suatu parametrisasi daerah itu, yakni bijeksi
$\phi:D\to G$ dari suatu subhimpunan $D$ dalam ruang Euklides. Namun, jelas
bahwa kita perlu memastikan hasil perhitungan tidak bergantung pada
parametrisasi yang dipilih. Walaupun teori ini mungkin saja didasarkan pada
hasil-hasil yang menunjukkan ketakbergantungan tersebut secara langsung,
hal itu menurut saya tidak sungguh mencerminkan intuisi yang mendasarinya,
yakni bahwa luas permukaan-permukaan sederhana, setidaknya, merupakan
sesuatu yang intrinsik bagi geometrinya. Karena itu, saya tidak melihat
alternatif yang dapat diterima selain teori `ukuran berdimensi $r$' yang
dapat diuraikan dengan istilah yang sepenuhnya geometris. Inilah pokok
\S264, tempat saya memberikan definisi dan sifat-sifat paling mendasar dari
ukuran Hausdorff berdimensi $r$ dalam ruang Euklides. Setelah dasar ini
tersedia, kita mendapati bahwa teknik-teknik \S\S261-263 cukup untuk
menghubungkannya dengan perhitungan melalui parametrisasi, sebagaimana saya
kerjakan dalam \S265.

Bab ini diakhiri dengan uraian singkat tentang ketaksamaan Brunn--Minkowski
(266C), yang merupakan alat penting bagi teori ukuran geometris himpunan
konveks.

\discrpage
